%!TEX TS-program = pdflatex
%!TEX encoding = UTF-8 Unicode
% LuxSleek-CV 1.1 LaTeX template
%
% This template has been adapted from the version published at:
% https://www.overleaf.com/latex/templates/luxsleek-cv/qbvbqmrzxwyj
%
% Author:
% Andreï V. Kostyrka, University of Luxembourg
%
% Adapted to be a Quarto template by Mitchell O'Hara-Wild
%
% Template license:
% CC BY 4.0 (https://creativecommons.org/licenses/by/4.0/), see LICENSE.md
%
% Class options, colours, the tinted-column machinery and Pandoc's own
% preamble -- everything before \begin{document}. Held in a partial, the way
% every other quarto-vitae template does, so a second format sharing this
% design (a cover letter, say) could not drift out of step with it.
%
% Upstream ships this as a single self-contained main.tex with zero page
% margins and two hand-sized minipages standing in for columns (see the note
% at the top of luxsleek.tex for why that does not survive a CV longer than one
% page). Everything here is the same document dressed instead in `paracol`,
% following the precedent already worked out for latexcv's sidebar designs.
%-------------------------------------------------------------------------------
% CONFIGURATIONS
%-------------------------------------------------------------------------------
$passoptions.latex()$
%
% A4 by default; upstream's own choice. Use 'letterpaper' for US letter.
$doc-class.tex()$

%-------------------------------------------------------------------------------
%	ENCODING
%-------------------------------------------------------------------------------
% Upstream is a pdflatex document; `iftex` guards these so the template still
% builds if a document overrides `pdf-engine` to xelatex or lualatex, both of
% which provide UTF-8 input and modern font handling natively.
\usepackage{iftex}
\ifPDFTeX
  \usepackage[T1]{fontenc}
  \usepackage[utf8]{inputenc}
\fi

%-------------------------------------------------------------------------------
%	LOGIC AND LAYOUT
%-------------------------------------------------------------------------------
% Two columns that flow independently and can be tinted the whole way down --
% see the "SIDEBAR" section of luxsleek.tex for what that buys over upstream's
% pair of minipages.
\usepackage{paracol}

% Page margins, set by luxsleek.tex (or overridden by the document's own
% `geometry:`). Upstream sets its own margins to zero and places everything by
% hand; paracol needs real margins to fill and tint instead, see luxsleek.tex.
\usepackage{geometry}

% Upstream's own list spacing, applied to any bullet list written in ordinary
% Markdown prose. The entries styles below do not use itemize themselves --
% `why` is set as a run of \smaller paragraphs, upstream's own device, not a
% true list -- so this only ever reaches a list a user writes by hand outside
% the entries system.
\usepackage{enumitem}
\setlist{parsep=0pt, topsep=0pt, partopsep=1pt, itemsep=1pt, leftmargin=6mm}

% Icons for the sidebar's contact block (\faIcon), exactly as the other
% quarto-vitae templates use fontawesome5 rather than the older v4 package.
\usepackage{fontawesome5}

\usepackage{graphicx}

% xcolor, for \definecolor/\textcolor/\colorbox -- fontawesome5 above happens
% to load it as a dependency, but it is asked for outright here rather than
% relied upon, since it is this file's own \definecolor{cvaccent}{...} below
% that needs it, not fontawesome5's.
\usepackage{xcolor}

% Honours `linestretch`, as every other Quarto pdf format does. Not loaded by
% common.latex() below, and upstream has no opinion on line spacing to
% inherit, so it is loaded outright here rather than only when linestretch is
% actually set in the front matter.
\usepackage{setspace}

%-------------------------------------------------------------------------------
%	FONT
%-------------------------------------------------------------------------------
% FiraSans is upstream's own choice, with the sans face made the document
% default, exactly as latexcv's designs default to Raleway. Skipped when the
% document sets `fontfamily`, so that option reaches font-settings.latex()
% below undisturbed instead of being fought over by two packages.
$if(fontfamily)$
$else$
\ifPDFTeX
  \usepackage{FiraSans}
\fi
\renewcommand*\familydefault{\sfdefault}
$endif$

% Upstream's own microtype configuration: `tracking`/`letterspace` are what
% version 1.1 added "for prettier lower caps" -- \headleft, \headright and
% every entry title here are set in \textsc, and letterspacing is what keeps
% small caps from reading as cramped at this weight.
%
% Passed as pending options rather than loaded outright: font-settings.latex(),
% called below via common.latex(), loads microtype itself if it is installed
% (`\IfFileExists{microtype.sty}{\usepackage[...]{microtype}}`), and a second
% `\usepackage{microtype}` here would be a package-option clash the moment the
% two disagree. \PassOptionsToPackage queues these for whichever \usepackage
% call loads the package first, so upstream's tuning survives without this
% template racing Pandoc's own template to `\usepackage{microtype}`.
\PassOptionsToPackage{stretch=25, shrink=25, tracking=true, letterspace=30}{microtype}

%-------------------------------------------------------------------------------
%	COLOUR
%-------------------------------------------------------------------------------
% Upstream's single `cvaccent`, used both as the sidebar's fill and as the
% colour \headright sets a main-column heading in -- one identity colour doing
% both jobs, exactly as upstream defines it. `headcolor` is the name every
% other quarto-vitae template uses for this, so it is read here rather than
% inventing a synonym; upstream's own near-black is kept as the default.
$if(headcolor)$
\definecolor{cvaccent}{HTML}{$headcolor$}
$else$
\definecolor{cvaccent}{HTML}{1A1A1A}
$endif$

%-------------------------------------------------------------------------------
%	PANDOC SUPPORT
%-------------------------------------------------------------------------------
% common.latex is required in every custom Quarto LaTeX template, and calls
% pandoc.tex, which in turn brings in syntax highlighting, tables.tex,
% graphics.tex, citations.tex, babel-lang.tex, tightlist.tex, biblio-config.tex
% and header-includes. Do not hand-copy any of those: they are Pandoc's to
% maintain, and a stale local copy silently drifts out of step with the reader.
%
% Upstream loads hyperref itself with `hidelinks`; common.latex() loads it too,
% so `hidelinks` is passed as an option here rather than repeating the
% \usepackage, which would be a package-clash error.
\PassOptionsToPackage{hidelinks}{hyperref}
$common.latex()$
$after-header-includes.latex()$
$hypersetup.latex()$

% Honours mainfont/mainfontoptions and their *fallback variants for a document
% rendered with xelatex/lualatex, the same way every other Quarto pdf format
% does. fonts.latex is NOT called: it would load lmodern over the font stack
% chosen above.
$font-settings.latex()$
